\documentclass[exercice,a4]{thedocument}%
\startdatakeys
    \source{laclassecooperativedestef.com}
    \BO{2026}
    \cycle{3}
    \competencesDuSocle{RE4,CO2}
    \objectifsApprentissage{C3C1ab,C3C1aj,C3C1ad,C3C2ac,C3C2ab,C3C1ae}
    \motsCles{perpendicularite, demi-cercles, triangles}
\enddatakeys
\begin{document}%
\noindent\textbf{Programme de construction 17}\par\smallskip\noindent
\begin{itemize}
    \item Tracer les droites $(\mathspoint{d})$ et $(\mathspoint{e})$
    perpendiculaires entre elles. Elles se coupent en $\mathspoint{A}$.
    \item Tracer le segment $[\mathspoint{BC}]$ tel que
    $\mathspoint{BC}=4\text{ cm}$ et que $\mathspoint{A}$ soit le milieu
    de $[\mathspoint{BC}]$.
    \item Tracer les demi-cercles $\mathcal{C}$ et $\mathcal{D}$,
    respectivement de centre $\mathspoint{B}$ et de centre
    $\mathspoint{C}$, de rayon $2\text{ cm}$.
    \item Le cercle $\mathcal{C}$ coupe la droite $(\mathspoint{d})$ en
    $\mathspoint{E}$ et le cercle $\mathcal{D}$ coupe $(\mathspoint{d})$
    en $\mathspoint{F}$.
    \item À l'opposé des demi-cercles, placer le point $\mathspoint{G}$ sur
    la droite $(\mathspoint{e})$ sachant que $\mathspoint{AG}=3\text{ cm}$.
    \item Tracer les segments $[\mathspoint{FG}]$ et $[\mathspoint{DG}]$.
\end{itemize}
\end{document}%
